% Fragment: §3-§4 of Bounded Deontics. Preamble, macros, theorem envs live in bd-draft.tex.

\section{The object level: a contract as a transition system}
\label{sec:object}

In L4, the ``letter'' of a contract is a state machine, and its primitive is deontically
\emph{neutral}. The core form names a party, an action, and the two branches that follow
from doing or not doing it:
\begin{verbatim}
    PARTY p DO a
      HENCE outcome_if_done
      LEST  outcome_if_not
\end{verbatim}
No moral weight is attached: \texttt{DO} merely records a choice point and its two
continuations. The familiar deontic operators are \emph{sugar} over this primitive, with
conventional default terminals---\texttt{MUST} fulfils on the action and breaches on the
deadline, \texttt{MAY} fulfils on either branch, \textsc{shant} breaches on the action:
\begin{verbatim}
    PARTY tenant MUST return_deposit WITHIN 30 days
      HENCE fulfilled    -- the action was taken
      LEST  breach       -- the deadline lapsed
\end{verbatim}
Crucially, \texttt{breach} and \texttt{fulfilled} are \emph{neutral terminals}, not
verdicts. They are Meyer's violation atom \src{Meyer 1988}: a marker that a designated
branch was taken, whose moral valence is supplied only by the enclosing context. An inner
contract's ``breach'' may be perfectly acceptable to the outer one---it records merely
``this path was not taken''---exactly as a Unix process returning a non-zero exit code
states a fact that the calling script is free to interpret. Contracts thus compose like
processes with exit codes, after the fashion of Hvitved's contract calculus
\src{Hvitved 2012}, with \texttt{HENCE}/\texttt{LEST} as the wiring (and, read as
synchronisation, the point of contact with process algebra \src{Hoare 1985}).

This object level is, by design, \emph{goal-free and preference-free}. It is precisely the
labelled transition system $T$ of Definition~\ref{def:dom}: states, actions, and the
branches between them, and nothing about which states are good. Goals and orderings enter
only at the assertion level, in Sections~\ref{sec:nec}--\ref{sec:boundary}. Keeping them
out of the letter is what lets one contract answer many different normative questions---and
is exactly the discipline that a goal-\emph{indexed} obligation gives up.

\section{Sorting constitutive from regulative: which ``musts'' are deontic}
\label{sec:sorter}

Before a ``must'' can be read as an obligation it must be sorted, because ordinary language
uses the word for two different things. Searle's distinction is the one we need
\src{Searle 1969; Searle 1995}. A \emph{regulative} rule prescribes behaviour---``directors
must vote within seven days''---and can be \emph{breached}: missing the deadline is a
violation with consequences. A \emph{constitutive} rule defines an institutional fact---
``all directors must vote yes \emph{for the resolution to pass}''---and cannot be breached
at all. A director who votes no violates nothing; the resolution simply does not pass. The
``must'' here is a \emph{must-be}, a necessary condition, not a duty.

The two readings are easy to conflate, because both wear the word ``must'' and both, at the
object level, are mere choice points and branches. The two-level architecture separates
them at the assertion level: a constitutive ``must $K$ for $J$'' is the bare statement that
$J$ is unreachable without $K$---there is no violation terminal on the no-$K$ branch, only
the non-occurrence of $J$---whereas a regulative ``must'' has a reachable \texttt{breach}
state that some party disvalues.

This is more than bookkeeping; it locates the seam of the paper. A constitutive must-be is
\emph{exactly} the party-neutral dominator relation of Section~\ref{sec:nec}:
$K\in\dom(J)$, necessity with no valence attached. A regulative must-\emph{do} is that same
relation \emph{plus} an ordering that makes failure a valued breach---the construction of
Section~\ref{sec:boundary}. The constitutive/regulative distinction is therefore not a
separate mechanism but the boundary between $\nect$ taken alone and $\nect$ composed with a
preference ordering. The sorter fixes the domain of every later quantifier: only regulative
musts carry obligation, and only they are read against the two orderings of
Section~\ref{sec:holmeshart}. Constitutive necessity needs no ``or else what?''---its ``or
else'' is just that the institutional fact fails to obtain.
